\begin{helper}[Replacement excursion and final-departure order]
Let $\beta,v,z,P_\delta$ be any replacement-coloring data with $v\in Z$,
with $\beta$ missing at $v$ under both colorings, free for $T_\delta$,
distinct from $\delta$ and $\gamma$, and with $P_\delta$ a full simple
$(\gamma,\beta)$ path from $v$ to $z$ outside $X_\delta$.  There is a
$\beta$-colored edge $h$ under $\varphi_\delta$ with exact ends
$u^1_\gamma,u^2_\gamma$.  In one of its two orientations,
$P_\delta$ contains the five consecutive nodes
\[
 f^1_\gamma,u^1_\gamma,h,u^2_\gamma,f^2_\gamma.
\]
This excursion occurs strictly before the last $X_\delta$-vertex; that
last vertex is $v^0_\delta$, the immediately following edge is
$f^0_\delta$, and all later vertex-nodes lie outside $X_\delta$.
Moreover neither $u^1_\delta$ nor $u^2_\delta$ is a Scheide exit of the
replacement tree.  The assertion retains exact finite indices for the
five consecutive nodes and the final vertex-edge pair.
\end{helper}
\begin{proof}
Because $\beta$ is missing under the original coloring at a vertex of
$Z$ and differs from $\gamma$,
\noderef{ScheideStructuralInternalPair} gives an original
$\beta$-edge $h$ internal to $U_\gamma$.  It is incident with $V(T)$,
is not $f^0_\delta$ because their original colors differ, and is therefore
fixed by \noderef{ScheideDefectiveSwitchEffect}.  Thus it still has color
$\beta$ under $\varphi_\delta$.  The exact configuration in
\noderef{ScheideSymmetricReplacementBoundaryConfiguration} makes both
deleted block edges $\gamma$-colored boundary identities, and
\noderef{ScheideDefectiveChainBoundaryEdges} puts both on the full path.
Properness of the replacement
coloring, looplessness, and
\noderef{TwoColorIncidenceDegreeAtMostTwo} now force the displayed five
nodes to be consecutive, in one orientation or the other, exactly as in
the original excursion argument.

We next exclude the two vertices of $U_\delta$ from the replacement exit
set.  By \noderef{ScheideGlobalMaximumMissingCounts}, each of them misses
at least one color.  Every color used on $T_\delta$ has a missing witness
in $Z$: retained base and block colors have their base witnesses from
\noderef{ScheideStructuralConfiguration}; those witnesses survive the
switch by \noderef{ScheideDefectiveSwitchEffect}, since distinct block
colors are injective, nonfirst base colors avoid all block colors, and
the defective color $\delta$ is unused by
\noderef{ScheideGlobalMaximumUnused}; finally, the appended color
$\delta$ is missing at $v^0_\delta\in Z$ by the switch-effect missing-set
formula.  The edge classification in \noderef{ScheideReplacementData}
shows that these retained edges and the appended pair exhaust the
displayed edges of $T_\delta$.  Elementariness from
\noderef{ScheideGlobalMaximumOrder} therefore says that no used color can
be missing at either vertex of $U_\delta$.  Any color missing there is
unused and hence free.  But \noderef{ScheideExitNoFree} says that an exit
of a globally maximum triple misses no free color.  Thus neither vertex
of $U_\delta$ is an exit.

The finite path starts in $X_\delta$ and ends outside it, so choose its
last vertex-node $w$ in $X_\delta$.  By the definition
\noderef{ScheideExitVertex}, $w$ is a replacement exit.  The preceding
paragraph and the exact equality $X_\delta=Z\cup U_\delta$ put $w$ in
$Z$.  Its next edge is a boundary edge.  Closedness excludes color
$\beta$, which is missing on $X_\delta$, so this edge has color $\gamma$.
The exact departure classification in
\noderef{ScheideSymmetricReplacementBoundaryConfiguration} makes it one
of $f^0_\delta,f^1_\gamma,f^2_\gamma$.

It cannot be either block edge, because the forced consecutive segment
then returns through the other block edge to a later vertex of $Z$.  It
is therefore $f^0_\delta$, whose exact tree endpoint is
$v^0_\delta$ by \noderef{ScheideDefectiveSetup}.  The returning block
excursion necessarily precedes this last departure.  Reading these
occurrences in the finite path order supplies the asserted consecutive
index equalities, strict order, final adjacent pair, and the condition on
all later vertex-nodes.
\end{proof}
